% Generated draft: requires mathematical and source review.
\documentclass[11pt]{article}
\usepackage[margin=25mm]{geometry}
\usepackage{fontspec}
\setmainfont{Latin Modern Roman}
\usepackage{amsmath,amssymb,mathtools,unicode-math}
\setmathfont{Latin Modern Math}
\usepackage{enumitem,needspace,xurl,hyperref}
\setlength{\parindent}{0pt}
\setlength{\parskip}{5pt}
\setlength{\emergencystretch}{3em}
\newcommand{\N}{\mathbb{N}}
\newcommand{\Z}{\mathbb{Z}}
\newcommand{\Q}{\mathbb{Q}}
\newcommand{\R}{\mathbb{R}}
\newcommand{\C}{\mathbb{C}}
\newcommand{\F}{\mathbb{F}}
\newcommand{\A}{\mathbb{A}}
\newcommand{\PP}{\mathbb{P}}
\newcommand{\E}{\mathbb{E}}
\newcommand{\eps}{\varepsilon}
\newcommand{\dd}{\,\mathrm{d}}
\begin{document}
\section*{UnsolvedMath ID: 10800009}
\textbf{Source problem code:} catalogue entry
\subsection{Definitions and mathematical statement}
Let \(E_d\subseteq\R^d\) be the \emph{essential ramification set} for the real-polynomial family intended in the source, and put \(X_d=\R^d\setminus E_d\). Vassiliev asks whether \(X_d\) is a \(K(\pi,1)\)-space. The defining polynomial family, coefficient coordinates, permitted \(d\), and definition of essential ramification are absent from this excerpt; \(E_d\) is a reference to that specific missing set, not the ordinary discriminant by default.

A path-connected space of CW type is a \(K(\pi,1)\) when its fundamental group is \(\pi\) and
\[\pi_j(X_d,x)=0\qquad(j\geq2),\]
with the usual local hypotheses giving the equivalent condition that its universal cover is contractible. If \(X_d\) is disconnected, the assertion must either specify a component or mean that every component is aspherical; a single connected Eilenberg--Mac Lane space and a disconnected union of such components are different conventions. The excerpt does not resolve this point. Thus a precise formulation, once \(E_d\) and the component convention are fixed, asks for vanishing of all higher homotopy groups, not merely a computation of \(\pi_1\) or homology.

The supplied context distinguishes this essential-ramification complement from classical discriminant/bifurcation complements, including the simple-singularity discriminant examples. Results about those other complements do not identify or settle this space without a comparison.
\subsection{Short English statement}
Is the complement of the specified essential ramification locus aspherical? The locus, degree range and connected-component convention are not defined in the short record and cannot be replaced by an ordinary polynomial discriminant.
\subsection{Sources}
\begin{itemize}
\item Local UnsolvedMath registry, permanent ID \(10800009\), exact source code \texttt{catalogue entry}.
\item Victor Vassiliev, \emph{A Few Problems on Monodromy and Discriminants}, Problem 3: \url{https://amj.math.stonybrook.edu/html-articles/Files-2015-2024/15-11/}. Registry access 2026-07-29; html-structure extraction.
\item V. A. Vassiliev, \emph{On topological invariants of real algebraic functions}, \emph{Functional Analysis and Its Applications} 45:3 (2011), 163--172; identified by the summary as source reference [essent].
\item The summary mentions Looijenga (1974) for simple-singularity discriminant complements, without a full entry, as context about a different space.
\end{itemize}
\end{document}
